\documentclass[DIN, pagenumber=false, fontsize=11pt, parskip=half]{scrartcl}
\usepackage[UTF8]{ctex}  % 中文支持
% \usepackage{ngerman}
\usepackage[utf8]{inputenc}
\usepackage[T1]{fontenc}
\usepackage{textcomp}
\usepackage{graphicx}

\usepackage{xcolor} % Required for custom red

\usepackage{geometry} % Custom margins

\usepackage{tocloft} % Modifying the table of contents

% for matlab code
% bw = blackwhite - optimized for print, otherwise source is colored
\usepackage[framed,numbered,bw]{mcode}
\usepackage[most]{tcolorbox}
% for other code
%\usepackage{listings}

\setlength{\parindent}{0em}

% set section in CM
\setkomafont{section}{\normalfont\bfseries\Large}

\newcommand{\mytitle}[1]{{\noindent\Large\textbf{#1}}}
\newcommand{\mysection}[1]{\textbf{\section*{#1}}}
\newcommand{\mysubsection}[2]{\romannumeral #1) #2}

\newtcolorbox{parchmentquote}{
    enhanced,
    colback=orange!5!white,   % 浅米黄色背景
    colframe=orange!50!black,
    boxrule=0.5pt,
    arc=2pt,                  % 细微圆角
    boxsep=5pt,
    before skip=10pt,
    after skip=10pt,
    fontupper=\itshape,
    borderline west={2pt}{0pt}{orange!80!black} % 加厚的左装饰线
}
%===================================
\begin{document}

\noindent\textbf{课后总结} \hfill \textbf{online course}\\
primary school \hfill Liu\\

\mytitle{乘法简单理解 \hfill \today}
\newline
\begin{parchmentquote}
  能够熟练乘法与加法的互相转换

\end{parchmentquote}


%===================================
\section{什么是乘法？}

乘法是一种\textbf{求几个相同加数和的简便运算}。

例如：

\[
5+5+5
\]

因为有 $3$ 个 $5$ 相加，

所以可以写成：

\[
3\times5=15
\]

读作：

\[
3\text{乘}5\text{等于}15
\]

意思是：

\[
\boxed{\text{3个5相加}}
\]
\section{计算}
\begin{parchmentquote}
 不一定非得用乘法算式才能得到结果，多思考乘法的由来

\end{parchmentquote}
\begin{enumerate}
    \item $156 \times 3 = $
    \item $567 \times 7 - 567 \times 5 = $
    \item $120 \times 5 = $

\end{enumerate}

\end{document}